% Источники: exam/materials/моделирование.pdf, раздел 17; exam/materials/прошлогодние ответы/Modelirovanie_1_-3.pdf, вопрос 20; https://github.com/HumsterProgrammer/iu9-notes/blob/main/8sem/Моделирование/Моделирование_лекция-06_03-03-26.md.
\section{Определение марковского случайного процесса. Моделирование марковских случайных процессов.}

\subsection*{Определение}

Случайный процесс $X(t)$ называется \textbf{марковским}, если для любых $t_1 < t_2 < \cdots < t_n < t$ выполняется условие:
\[
  P\bigl(X(t) \leqslant x \mid X(t_1), X(t_2), \ldots, X(t_n)\bigr)
  = P\bigl(X(t) \leqslant x \mid X(t_n)\bigr).
\]
Иными словами, \textbf{будущее зависит только от настоящего состояния}, но не от прошлого («свойство отсутствия последействия»).

\subsection*{Марковская цепь (дискретное время)}

Множество состояний $S = \{s_1, \ldots, s_m\}$; переходы происходят в дискретные моменты.
\textbf{Матрица переходных вероятностей} $P = (p_{ij})$, $p_{ij} = P(X_{n+1}=s_j\mid X_n=s_i)$, $\sum_j p_{ij}=1$.

Моделирование: в состоянии $s_i$ генерируется $r \sim U(0,1)$; по кумулятивным вероятностям строки $i$ выбирается следующее состояние.

\textbf{Стационарное распределение} $\boldsymbol{\pi}$: $\boldsymbol{\pi}P = \boldsymbol{\pi}$, $\sum_i \pi_i = 1$.

\subsection*{Непрерывная марковская цепь}

Переходы происходят в случайные моменты времени. Задаётся матрицей интенсивностей $\Lambda = (\lambda_{ij})$, где $\lambda_{ij} \geqslant 0$ при $i \neq j$ — интенсивность перехода $s_i \to s_j$.

\textbf{Уравнения Колмогорова} для вероятностей состояний $p_i(t) = P(X(t)=s_i)$:
\[
  \frac{dp_i}{dt} = \sum_{j \neq i} p_j \lambda_{ji} - p_i \sum_{j \neq i} \lambda_{ij}.
\]

Моделирование: в состоянии $s_i$ время до следующего перехода $\tau \sim \mathrm{Exp}\!\left(\sum_{j\neq i}\lambda_{ij}\right)$; тип перехода выбирается с вероятностью $\lambda_{ij}/\sum_{k\neq i}\lambda_{ik}$.

\clearpage
